\documentclass[12pt]{ximera}
\typeout{Start loading xmPreamble.tex}%

% Add here extra macro's that are loaded automatically by all documents of claas 'ximera' or 'xourse' in this repo

%%
%%  Example:
%%
% \newcommand{\R}{\mathbb{R}}
\usepackage{tikz}
\usetikzlibrary{positioning, arrows.meta}
\usepackage{multicol}
\usepackage{enumitem}
\usepackage{caption}
\usepackage{amsmath}
\usepackage{amsthm}
\usepackage{fontenc}

%% ---------------------------------------------------------------
%% Shared worksheet macros.
%%
%% These came from the Summer 2026 worksheet set, where each worksheet carried its
%% own copy. They have to live here instead: a xourse inlines every activity into a
%% single document, so duplicate \newcommand definitions across activities collide,
%% and the xourse gobbles \input so a per-topic macro file never loads.
%%
%% They use \providecommand, NOT \newcommand. ximera.cls loads this file automatically,
%% and every activity also carries an explicit \typeout{Start loading xmPreamble.tex}%

% Add here extra macro's that are loaded automatically by all documents of claas 'ximera' or 'xourse' in this repo

%%
%%  Example:
%%
% \newcommand{\R}{\mathbb{R}}
\usepackage{tikz}
\usetikzlibrary{positioning, arrows.meta}
\usepackage{multicol}
\usepackage{enumitem}
\usepackage{caption}
\usepackage{amsmath}
\usepackage{amsthm}
\usepackage{fontenc}

%% ---------------------------------------------------------------
%% Shared worksheet macros.
%%
%% These came from the Summer 2026 worksheet set, where each worksheet carried its
%% own copy. They have to live here instead: a xourse inlines every activity into a
%% single document, so duplicate \newcommand definitions across activities collide,
%% and the xourse gobbles \input so a per-topic macro file never loads.
%%
%% They use \providecommand, NOT \newcommand. ximera.cls loads this file automatically,
%% and every activity also carries an explicit \typeout{Start loading xmPreamble.tex}%

% Add here extra macro's that are loaded automatically by all documents of claas 'ximera' or 'xourse' in this repo

%%
%%  Example:
%%
% \newcommand{\R}{\mathbb{R}}
\usepackage{tikz}
\usetikzlibrary{positioning, arrows.meta}
\usepackage{multicol}
\usepackage{enumitem}
\usepackage{caption}
\usepackage{amsmath}
\usepackage{amsthm}
\usepackage{fontenc}

%% ---------------------------------------------------------------
%% Shared worksheet macros.
%%
%% These came from the Summer 2026 worksheet set, where each worksheet carried its
%% own copy. They have to live here instead: a xourse inlines every activity into a
%% single document, so duplicate \newcommand definitions across activities collide,
%% and the xourse gobbles \input so a per-topic macro file never loads.
%%
%% They use \providecommand, NOT \newcommand. ximera.cls loads this file automatically,
%% and every activity also carries an explicit \input{../xmPreamble}, so this file is
%% read twice. That was harmless while it held only \usepackage lines (LaTeX skips a
%% package it has already loaded) but a second \newcommand is a hard error.
%%
%%   \blank[width]        fill-in-the-blank rule (default 2.5cm)
%%   \showwork            "show and explain your work" reminder
%%   \showworkline        ... plus which schedule line was used
%%   \showworkyear        ... plus which year's schedule was used
%%   \schedhead{...}      centred bold caption above a tiered-rate table
%%   \samesquares[scale]{left}{right}   two equal-area squares, labelled
%%   \samecubes[scale]{left}{right}     two equal-volume cubes, labelled
%% ---------------------------------------------------------------

\providecommand{\blank}[1][2.5cm]{\underline{\hspace{#1}}}
\providecommand{\showwork}{\textbf{REMEMBER TO SHOW AND EXPLAIN YOUR WORK.}}
\providecommand{\showworkline}{\textbf{REMEMBER TO SHOW AND EXPLAIN YOUR WORK.
  Explanation should include how and why you know which line of the
  schedule was used to calculate this amount.}}
\providecommand{\showworkyear}{\begin{sloppypar}\textbf{REMEMBER TO SHOW AND
  EXPLAIN YOUR WORK. Explanation should include how and why you know which
  line of the year's tax schedule was used to calculate this tax
  amount.}\end{sloppypar}}
\providecommand{\schedhead}[1]{\begin{center}\textbf{#1}\end{center}}

\providecommand{\samesquares}[3][1]{%
\begin{center}
{\footnotesize These squares are the \textbf{same size}.}

\vspace{1.5mm}
\begin{tikzpicture}[scale=#1,every node/.style={scale=#1}]
  \draw (0,0) rectangle (2.4,2.4);
  \node[anchor=north west] at (0.15,2.25) {Area:};
  \node[anchor=east]  at (-0.15,1.2) {#2};
  \node[anchor=north] at (1.2,-0.15) {#2};
  \begin{scope}[xshift=4.6cm]
    \draw (0,0) rectangle (2.4,2.4);
    \node[anchor=north west] at (0.15,2.25) {Area:};
    \node[anchor=east]  at (-0.15,1.2) {#3};
    \node[anchor=north] at (1.2,-0.15) {#3};
  \end{scope}
\end{tikzpicture}
\end{center}}

\providecommand{\samecubes}[3][1]{%
\begin{center}
{\footnotesize These cubes are the \textbf{same size}.}

\vspace{1.5mm}
\begin{tikzpicture}[scale=#1,every node/.style={scale=#1}]
  \foreach \dx in {0,5.2} {
    \begin{scope}[xshift=\dx cm]
      \draw (0,0) rectangle (2.0,2.0);
      \draw (0,2.0) -- (0.65,2.6) -- (2.65,2.6) -- (2.0,2.0);
      \fill[black!12] (2.0,0) -- (2.65,0.6) -- (2.65,2.6) -- (2.0,2.0) -- cycle;
      \draw (2.0,0) -- (2.65,0.6) -- (2.65,2.6) -- (2.0,2.0);
      \node[anchor=north west] at (0.1,1.9) {Volume:};
    \end{scope}
  }
  \node[anchor=east]  at (-0.15,1.0) {#2};
  \node[anchor=north] at (1.0,-0.15) {#2};
  \node[anchor=west]  at (2.25,0.4) {#2};
  \node[anchor=east]  at (5.05,1.0) {#3};
  \node[anchor=north] at (6.2,-0.15) {#3};
  \node[anchor=west]  at (7.45,0.4) {#3};
\end{tikzpicture}
\end{center}}

%% NOTE: do NOT add \usepackage to individual activity .tex files.
%% When a xourse inlines an activity it gobbles \documentclass and \input, but a bare
%% \usepackage still executes -- after \begin{document} -- and errors out.
%% All shared packages and macros belong here.
%%
%% ximera.cls already loads: amsmath, amssymb, amsthm, cancel, enumitem, graphicx,
%% hyperref, listings, pgfplots, tikz, url, xcolor. No need to re-load those.
, so this file is
%% read twice. That was harmless while it held only \usepackage lines (LaTeX skips a
%% package it has already loaded) but a second \newcommand is a hard error.
%%
%%   \blank[width]        fill-in-the-blank rule (default 2.5cm)
%%   \showwork            "show and explain your work" reminder
%%   \showworkline        ... plus which schedule line was used
%%   \showworkyear        ... plus which year's schedule was used
%%   \schedhead{...}      centred bold caption above a tiered-rate table
%%   \samesquares[scale]{left}{right}   two equal-area squares, labelled
%%   \samecubes[scale]{left}{right}     two equal-volume cubes, labelled
%% ---------------------------------------------------------------

\providecommand{\blank}[1][2.5cm]{\underline{\hspace{#1}}}
\providecommand{\showwork}{\textbf{REMEMBER TO SHOW AND EXPLAIN YOUR WORK.}}
\providecommand{\showworkline}{\textbf{REMEMBER TO SHOW AND EXPLAIN YOUR WORK.
  Explanation should include how and why you know which line of the
  schedule was used to calculate this amount.}}
\providecommand{\showworkyear}{\begin{sloppypar}\textbf{REMEMBER TO SHOW AND
  EXPLAIN YOUR WORK. Explanation should include how and why you know which
  line of the year's tax schedule was used to calculate this tax
  amount.}\end{sloppypar}}
\providecommand{\schedhead}[1]{\begin{center}\textbf{#1}\end{center}}

\providecommand{\samesquares}[3][1]{%
\begin{center}
{\footnotesize These squares are the \textbf{same size}.}

\vspace{1.5mm}
\begin{tikzpicture}[scale=#1,every node/.style={scale=#1}]
  \draw (0,0) rectangle (2.4,2.4);
  \node[anchor=north west] at (0.15,2.25) {Area:};
  \node[anchor=east]  at (-0.15,1.2) {#2};
  \node[anchor=north] at (1.2,-0.15) {#2};
  \begin{scope}[xshift=4.6cm]
    \draw (0,0) rectangle (2.4,2.4);
    \node[anchor=north west] at (0.15,2.25) {Area:};
    \node[anchor=east]  at (-0.15,1.2) {#3};
    \node[anchor=north] at (1.2,-0.15) {#3};
  \end{scope}
\end{tikzpicture}
\end{center}}

\providecommand{\samecubes}[3][1]{%
\begin{center}
{\footnotesize These cubes are the \textbf{same size}.}

\vspace{1.5mm}
\begin{tikzpicture}[scale=#1,every node/.style={scale=#1}]
  \foreach \dx in {0,5.2} {
    \begin{scope}[xshift=\dx cm]
      \draw (0,0) rectangle (2.0,2.0);
      \draw (0,2.0) -- (0.65,2.6) -- (2.65,2.6) -- (2.0,2.0);
      \fill[black!12] (2.0,0) -- (2.65,0.6) -- (2.65,2.6) -- (2.0,2.0) -- cycle;
      \draw (2.0,0) -- (2.65,0.6) -- (2.65,2.6) -- (2.0,2.0);
      \node[anchor=north west] at (0.1,1.9) {Volume:};
    \end{scope}
  }
  \node[anchor=east]  at (-0.15,1.0) {#2};
  \node[anchor=north] at (1.0,-0.15) {#2};
  \node[anchor=west]  at (2.25,0.4) {#2};
  \node[anchor=east]  at (5.05,1.0) {#3};
  \node[anchor=north] at (6.2,-0.15) {#3};
  \node[anchor=west]  at (7.45,0.4) {#3};
\end{tikzpicture}
\end{center}}

%% NOTE: do NOT add \usepackage to individual activity .tex files.
%% When a xourse inlines an activity it gobbles \documentclass and \input, but a bare
%% \usepackage still executes -- after \begin{document} -- and errors out.
%% All shared packages and macros belong here.
%%
%% ximera.cls already loads: amsmath, amssymb, amsthm, cancel, enumitem, graphicx,
%% hyperref, listings, pgfplots, tikz, url, xcolor. No need to re-load those.
, so this file is
%% read twice. That was harmless while it held only \usepackage lines (LaTeX skips a
%% package it has already loaded) but a second \newcommand is a hard error.
%%
%%   \blank[width]        fill-in-the-blank rule (default 2.5cm)
%%   \showwork            "show and explain your work" reminder
%%   \showworkline        ... plus which schedule line was used
%%   \showworkyear        ... plus which year's schedule was used
%%   \schedhead{...}      centred bold caption above a tiered-rate table
%%   \samesquares[scale]{left}{right}   two equal-area squares, labelled
%%   \samecubes[scale]{left}{right}     two equal-volume cubes, labelled
%% ---------------------------------------------------------------

\providecommand{\blank}[1][2.5cm]{\underline{\hspace{#1}}}
\providecommand{\showwork}{\textbf{REMEMBER TO SHOW AND EXPLAIN YOUR WORK.}}
\providecommand{\showworkline}{\textbf{REMEMBER TO SHOW AND EXPLAIN YOUR WORK.
  Explanation should include how and why you know which line of the
  schedule was used to calculate this amount.}}
\providecommand{\showworkyear}{\begin{sloppypar}\textbf{REMEMBER TO SHOW AND
  EXPLAIN YOUR WORK. Explanation should include how and why you know which
  line of the year's tax schedule was used to calculate this tax
  amount.}\end{sloppypar}}
\providecommand{\schedhead}[1]{\begin{center}\textbf{#1}\end{center}}

\providecommand{\samesquares}[3][1]{%
\begin{center}
{\footnotesize These squares are the \textbf{same size}.}

\vspace{1.5mm}
\begin{tikzpicture}[scale=#1,every node/.style={scale=#1}]
  \draw (0,0) rectangle (2.4,2.4);
  \node[anchor=north west] at (0.15,2.25) {Area:};
  \node[anchor=east]  at (-0.15,1.2) {#2};
  \node[anchor=north] at (1.2,-0.15) {#2};
  \begin{scope}[xshift=4.6cm]
    \draw (0,0) rectangle (2.4,2.4);
    \node[anchor=north west] at (0.15,2.25) {Area:};
    \node[anchor=east]  at (-0.15,1.2) {#3};
    \node[anchor=north] at (1.2,-0.15) {#3};
  \end{scope}
\end{tikzpicture}
\end{center}}

\providecommand{\samecubes}[3][1]{%
\begin{center}
{\footnotesize These cubes are the \textbf{same size}.}

\vspace{1.5mm}
\begin{tikzpicture}[scale=#1,every node/.style={scale=#1}]
  \foreach \dx in {0,5.2} {
    \begin{scope}[xshift=\dx cm]
      \draw (0,0) rectangle (2.0,2.0);
      \draw (0,2.0) -- (0.65,2.6) -- (2.65,2.6) -- (2.0,2.0);
      \fill[black!12] (2.0,0) -- (2.65,0.6) -- (2.65,2.6) -- (2.0,2.0) -- cycle;
      \draw (2.0,0) -- (2.65,0.6) -- (2.65,2.6) -- (2.0,2.0);
      \node[anchor=north west] at (0.1,1.9) {Volume:};
    \end{scope}
  }
  \node[anchor=east]  at (-0.15,1.0) {#2};
  \node[anchor=north] at (1.0,-0.15) {#2};
  \node[anchor=west]  at (2.25,0.4) {#2};
  \node[anchor=east]  at (5.05,1.0) {#3};
  \node[anchor=north] at (6.2,-0.15) {#3};
  \node[anchor=west]  at (7.45,0.4) {#3};
\end{tikzpicture}
\end{center}}

%% NOTE: do NOT add \usepackage to individual activity .tex files.
%% When a xourse inlines an activity it gobbles \documentclass and \input, but a bare
%% \usepackage still executes -- after \begin{document} -- and errors out.
%% All shared packages and macros belong here.
%%
%% ximera.cls already loads: amsmath, amssymb, amsthm, cancel, enumitem, graphicx,
%% hyperref, listings, pgfplots, tikz, url, xcolor. No need to re-load those.

\title{Exponential Decay: Rolling Dice Away}
\begin{document}
\begin{abstract}
A hands-on decay experiment --- roll a pile of dice, remove every die showing a ``1,''
repeat until none are left --- and then fit an exponential model to the class's pooled
data. Converting between the $y = a e^{bx}$ form that spreadsheets produce and the
$y = a R^x$ form that makes the context readable.
\end{abstract}
\maketitle

\section*{The Experiment}

Go to \url{https://www.randomlists.com/dice} to simulate dice rolling. Set the number of
sides to 6 and the total number of dice to roll to 25, then click ``return'' to simulate
a roll.

Roll the 25 dice and count how many came up as a ``1.'' Remove those dice from your
total. So if you rolled 25 dice and 7 came up ``1,'' then for the next roll you change
your total to $25 - 7 = 18$ dice.

If no dice roll a ``1,'' then none are removed --- but you still record that value again
in your table for that roll. Especially when few dice are left, you may have several
rolls in a row where the number does not change.

Continue rolling and recording how many dice remain until you have 0 dice left. For the
first 10 rolls, also calculate what percentage of dice remain from the previous roll.

\subsection*{Example data}

\begin{center}
\begin{tabular}{ccl}
\hline
Roll & Dice remaining & \% remaining from previous roll \\
\hline
0  & 25 & N/A \\
1  & 21 & $21/25 = 0.840 = 84.0\%$ \\
2  & 19 & $19/21 = 0.905 = 90.5\%$ \\
3  & 13 & $13/19 = 0.684 = 68.4\%$ \\
4  & 11 & $11/13 = 0.846 = 84.6\%$ \\
5  & 10 & $10/11 = 0.909 = 90.9\%$ \\
6  & 9  & $9/10 = 0.900 = 90.0\%$ \\
7  & 9  & $9/9 = 1.000 = 100\%$ \\
8  & 7  & $7/9 = 0.778 = 77.8\%$ \\
9  & 7  & $7/7 = 1.000 = 100\%$ \\
10 & 6  & $6/7 = 0.857 = 85.7\%$ \\
11 & 5  & \\
12 & 4  & \\
13 & 3  & \\
14 & 3  & \\
15 & 2  & \\
16 & 1  & \\
17 & 1  & \\
18 & 0  & \\
\hline
\end{tabular}
\end{center}

This example reached 0 dice at roll 18. Yours may take more or fewer rolls --- the two
rolls above where the count did not change (rolls 7 and 9) are exactly the kind of thing
that makes the length vary.

Once you have rolled away all your dice, enter your dice-remaining values for each roll
into your class's Google Sheet. While the class finishes, work on the following.

\section*{The Other Form: $y = a \cdot e^{bx}$}

The exponential functions we have worked with have all been written as $y = a \cdot R^x$.
But exponential models can also come in the form $y = a \cdot e^{bx}$, and this second
form is what programs most commonly produce.

That ``$e$'' is just a letter for a special irrational number. Just as
$\pi = 3.14159\ldots$, we have $e = 2.71828\ldots$ Your scientific calculator should have
a button for it.

\begin{problem}
Suppose we have an exponential model $y = 671 \cdot e^{0.3x}$. Rewrite this in the form
$y = a \cdot R^x$.

Using properties of exponents, $e^{bx} = (e^b)^x$. Calculate the $e^b$ value, rounded to
four decimal places.

$e^{0.3} = \answer[tolerance=0.001]{1.3499}$

\begin{freeResponse}
\end{freeResponse}

\begin{solution}
\[
e^{0.3} = 1.3499\ldots
\]
which means
\[
y = 671 \cdot e^{0.3x} = 671 \cdot (1.3499)^x.
\]
The whole conversion rests on the exponent rule $e^{bx} = (e^b)^x$ --- once you see that,
$R$ is just $e^b$.
\end{solution}
\end{problem}

\section*{Modelling the Class Data}

Once everyone has entered their data, your instructor will have the spreadsheet produce a
model. In the example data below, the class rolled 6-sided dice, actually started with
\textbf{129} dice, and the spreadsheet produced
\[
y = 136 \cdot e^{-0.181x}.
\]

\begin{problem}
What do the $x$-value and the $y$-value represent in this model?

\begin{freeResponse}
\end{freeResponse}

\begin{solution}
The $x$-values represent the number of rolls. The $y$-values represent the number of dice
remaining after $x$ rolls.
\end{solution}
\end{problem}

\begin{problem}
The starting number of dice should be when $x = \answer{0}$. Use this and the model to
determine the model's estimate for the starting number of dice.

$\answer{136}$ dice

\begin{solution}
The starting number is the count at roll 0, before any rolls have occurred --- hence
$x = 0$, since $x$ represents the number of rolls. Plugging in:
\[
y(0) = 136 \cdot e^{-0.181(0)} = 136 \cdot e^{0} = 136 \cdot 1 = 136.
\]
\end{solution}
\end{problem}

\begin{problem}
Why does the model's estimate of 136 not match the actual starting number of 129?

\begin{freeResponse}
\end{freeResponse}

\begin{solution}
Because the data is not perfectly exponential. The spreadsheet finds the exponential
curve that best fits \emph{all} the data points taken together, and that best-fit curve is
not required to pass exactly through any particular one of them --- including the starting
count.

Real dice rolls are random, so the class's actual counts wobble around the smooth
exponential shape rather than sitting on it. The model trades a little accuracy at roll 0
for a better overall fit across every roll.
\end{solution}
\end{problem}

\begin{problem}
Now plug $x = 8$ into the model. Explain what the $x$ and $y$ values mean in this context.

$y(8) \approx \answer[tolerance=1]{32}$ dice

\begin{freeResponse}
\end{freeResponse}

\begin{solution}
\[
y(8) = 136 \cdot e^{-0.181(8)} = 136 \cdot (0.2350\ldots) = 31.97 \approx 32.
\]
This means that when $x = 8$, $y = 32$: after 8 rolls, we should expect about 32 dice
remaining.
\end{solution}
\end{problem}

\section*{Rewriting the Model}

We can rewrite the model the program produced without the ``$e$,'' putting it in a more
recognisable form.

\begin{problem}
Using $e^{bx} = (e^b)^x$ and the numbers from our model, calculate the $e^b$ value,
rounded to four decimal places, and write the model in the form $y = a \cdot R^x$.

$e^{-0.181} = \answer[tolerance=0.001]{0.8344}$

\begin{freeResponse}
\end{freeResponse}

\begin{solution}
\[
e^{-0.181} = 0.8344\ldots
\]
so the model becomes
\[
y = 136 \cdot (0.8344)^x.
\]
\end{solution}
\end{problem}

\begin{problem}
What should this $R$-value be representing in this context? What percentage of dice are
lost with each roll?

$\answer[tolerance=0.2]{16.56}\%$ lost per roll

\begin{freeResponse}
\end{freeResponse}

\begin{solution}
The $R$-value tells us about the decay rate --- the percentage change in the number of
dice after each roll.

In our earlier examples we had exponential \emph{growth}, so $R$ showed by what percentage
values were increasing: an $R$ of 1.056 meant $1.056 = 1 + 0.056$, a $5.6\%$ increase per
year. Here we have exponential \emph{decrease}, so
\[
0.8344 = 1 - 0.1656 = 100\% - 16.56\%,
\]
meaning about a $16.56\%$ decrease per roll. We expect to lose about $16.56\%$ of our dice
with each roll.

Equivalently --- and often the more natural reading --- if we lose $16.56\%$, then
$83.44\%$ of the dice \emph{remain} after each roll. The $R$-value is directly the
fraction surviving.
\end{solution}
\end{problem}

\begin{problem}
Now plug $x = 8$ into the rewritten model. Does your value equal what you got from the
other form? Should they be equal? Why or why not?

$y(8) \approx \answer[tolerance=1]{32}$ dice

\begin{freeResponse}
\end{freeResponse}

\begin{solution}
\[
y(8) = 136 \cdot (0.8344)^8 = 136 \cdot (0.2350\ldots) = 31.95 \approx 32.
\]

The two answers round to the same value but are not identical past the decimal point:
31.97 from the $e$ form versus 31.95 from the $R$ form.

They \emph{should} be equal in principle, because both forms are the same underlying
model just written differently. The small discrepancy comes entirely from rounding $e^b$
to four decimal places when we wrote $R = 0.8344$. Carrying full precision would make
them match exactly.

This is worth noticing generally: rounding an intermediate value introduces error that
gets amplified when that value is raised to a power.
\end{solution}
\end{problem}

\end{document}
